\section{Fabrication}

All of the devices reported in this thesis were nominally made using the following recipe. The recipe was developed from a recipe developed in the Javad Shabani group at NYU.


\begin{enumerate}
    \item \textbf{Dice}
        \begin{enumerate}
            \item Spin AZ1815 photoresist at 3000rpm for 30s. 
            \item Bake for 2 min at 110$\degree$C
            \item Dice using Advanced Dicing Technologies (ADT) 7112 Wafer Dicer with Ni blade (ADT 00757-4240-890-200) at 10 krpm and 1mm/s cut rotation and speed, respectively.
            \item Standard solvent clean
            \begin{enumerate}
                \item Soak in NMP for 30 mins at 80\degree C.
                \item Soak in Acetone for 1 min.
                \item Soak in IPA for 1 min.
                \item Dry using \ce{N2} gun
                
            \end{enumerate}
            
        \end{enumerate}
    \item \textbf{Shallow trench etch} This step defines the width of the junction, the RF SQUID loop, alignment markers, and the gap of the gate line.
        \begin{enumerate}
            \item Spin PMMA 950 A4 at 4000 rpm (at 1000 rpm/s ramp up) for 60s.
            \item Bake at 180\degree C for 120 s.
            \item Expose in Raith EBPG 5150 with 20nA beam current and nominally 700 $\mu C/cm^2 $
            \item Develop in  1:3, DI water:IPA solution at 5\degree for 2 min. Dry using a \ce{N2} gun.
            \item Bake 80\degree C for 5 min. Helps prevent peeling during acid etch.
            \item Al acid etch:
            \begin{enumerate}
                \item Place beaker of Transene Aluminium Etchant Type D in water bath on hot plate and set temperature to 50\degree C. The water temperature should end up being between 29-32 \degree C
                \item Dip sample in etchant until color of large features changes. Wait an additional 3 seconds.
                \item Triple rinse in DI water and blow dry with \ce{N2} gun.
            \end{enumerate}
            \item Dip sample in Dilute Phosphoric Peroxide\footnote{I had very consistent etch results ordering this mixture from Transene} (DPP, Phosphoric acid(85\%): Hydrogen Peroxide(30\%): DI water in a volume ratio of 1:1:40) for 2 mins. Etches ~200nm of substrate.
            \item Triple rinse in DI water and blow dry with \ce{N2} gun.
            \item Dip sample back in Transene Aluminium Etchant Type D for about 5 seconds. This removes the overhung Al from the undercutting DPP etch.
            \item Triple rinse in DI water and blow dry with \ce{N2} gun.
            \item Standard solvent clean
            \begin{enumerate}
                \item Soak in NMP for 30 mins at 80\degree C.
                \item Soak in Acetone for 1 min.
                \item Soak in IPA for 1 min.
                \item Dry using \ce{N2} gun
                
            \end{enumerate}
            
        \end{enumerate}
    \item \textbf{Deep Trench Etch} This step defines the microwave components like the resonator, and feedline.
    \begin{enumerate}
        \item Spin AZ 1518 photo-resist at 4000 rpm (1000 rpm/s ramp up) for 90s.
        \item Bake at 100\degree C for 50 s.
        \item Expose in Heidelberg MLA 150 with 375 nm light at $65mJ/cm^2 $
        \item Develop in AZ 917 MIF for 85 s. Rinse with DI water and dry using a \ce{N2} gun.
        \item Al acid etch (If TMAH in developer didn't fully etch Al):
        \begin{enumerate}
            \item Place beaker of Transene Aluminium Etchant Type D in water bath on hot plate and set temperature to 50\degree C. The water temperature should end up being between 29-32 \degree C
            \item Dip sample in etchant until color of large features changes. Wait an additional 3 seconds.
            \item Triple rinse in DI water and blow dry with \ce{N2} gun.
        \end{enumerate}
        \item Dip in DPP for 5 min. Nominally etches ~500 nm into substrate.
        \item Triple rinse in DI water and blow dry with \ce{N2} gun.
        \item Standard solvent clean
        \begin{enumerate}
            \item Soak in NMP for 30 mins at 80\degree C.
            \item Soak in Acetone for 1 min.
            \item Soak in IPA for 1 min.
            \item Dry using \ce{N2} gun
        \end{enumerate}
    \end{enumerate}
    \item \textbf{Aluminum junction etch} This step defines the length of the junction.
    \begin{enumerate}
        \item Spin PMMA 950 A4 at 4000 rpm (at 1000 rpm/s ramp up) for 60 s.
        \item Bake at 180\degree C for 120 s.
        \item Expose in Raith EBPG 5150 with 20 nA beam current and nominally 700 $\mu C/cm^2 $
        \item Develop in  1:3, DI water:IPA solution at 5\degree for 2 min. Dry using a \ce{N2} gun.
        \item Bake 80\degree C for 5 min. Helps prevent peeling during acid etch.
        \item Al acid etch:
        \begin{enumerate}
            \item Place beaker of Transene Aluminium Etchant Type D in water bath on hot plate and set temperature to 50\degree C. The water temperature should end up being between 29-32 \degree C
            \item Dip sample in etchant until color of large features changes. Wait an additional 3 seconds.
            \item Triple rinse in DI water and blow dry with \ce{N2} gun.
        \end{enumerate}
        \item Standard solvent clean
        \begin{enumerate}
            \item Soak in NMP for 30 mins at 80\degree C.
            \item Soak in Acetone for 1 min.
            \item Soak in IPA for 1 min.
            \item Dry using \ce{N2} gun
        \end{enumerate}
    \end{enumerate}
    \item \textbf{\ce{HfO2} Dielectric (Lifted off)} This defines the dielectric to isolate gate from junction or RF SQUID in a defined area.
    \begin{enumerate}
        \item Spin PMMA 950 A4 at 4000 rpm (at 1000 rpm/s ramp up) for 60 s.
        \item Bake at 180\degree C for 120 s.
        \item Expose in Raith EBPG 5150 with 20 nA beam current and nominally 700 $\mu C/cm^2 $
        \item Develop in  1:3, DI water:IPA solution at 5\degree for 2 min. Dry using a \ce{N2} gun.
        \item Atomic layer deposition in Savannah S100 Atomic Layer Deposition Cambridge Nanotech.
        \begin{enumerate}
            \item Flow \ce{N2} at 20 sscm. Chamber temp is set to 120$\degree$C.
            \item Repeat the following for 300 cycles:
            \begin{enumerate}
                \item Pulse water valve for 0.1 s
                \item Wait 40 s.
                \item Pulse Hf precursor for 0.3 s.
                \item Wait 40 s.
            \end{enumerate}
        \end{enumerate}
        \item Solvent liftoff
        \begin{enumerate}
            \item Soak in NMP for 120 mins at 80\degree C.
            \item Sonnicate for 1 min in NMP.
            \item Soak in Acetone for 1 min.
            \item Soak in IPA for 1 min.
            \item Dry using \ce{N2} gun
        \end{enumerate}
    \end{enumerate}
    \item \textbf{Gate liftoff} This defines the gate.
    \begin{enumerate}
        \item Spin MMA EL13 at 4000 rpm (at 1000 rpm/s ramp up) for 60 s.
        \item Bake at 180\degree C for 120 s.
        \item Spin PMMA 950 A4 at 4000 rpm (at 1000 rpm/s ramp up) for 60 s.
        \item Bake at 180\degree C for 120 s.
        \item Expose in Raith EBPG 5150 with 20 nA beam current and nominally 700 $\mu C/cm^2 $
        \item Develop in  1:3, DI water:IPA solution at 5\degree for 2 min. Dry using a \ce{N2} gun.
        \item Sputter about 230 nm Nb at 0.8 nm/s and 6 mTorr Ar pressure.
        \item Solvent liftoff:
        \begin{enumerate}
            \item Soak in NMP for 120 mins at 80\degree C.
            \item While sample is still submerged in NMP. Spray to remove any film remaining on surface.
            \item Soak in Acetone for 1 min.
            \item Soak in IPA for 1 min.
            \item Dry using \ce{N2} gun
        \end{enumerate}
    \end{enumerate}
    
\end{enumerate}

